\section{Results and Discussion}
\label{sec:results}

% Rubric (Analysis & Interpretation): "Results analyzed thoroughly with
% appropriate reasoning. Limitations and alternative explanations discussed."
% Guidelines:
%   - Explore and compare multiple configurations (hyperparameters, hidden
%     states, structures); explain differences in performance and stability.
%   - Quantitative results (likelihoods, accuracies, ...).
%   - Qualitative insights: what the model captures or misses; do learned
%     CPTs / latent states match intuition?
%   - Convergence and scalability.
%
% Tip: every row in a results table should come from a results/<run>/ folder
% and a linked Experiment issue.

\subsection{Experimental Setup}
% Metrics, splits, configurations compared, seeds / restarts.
\todo{Setup.}

\subsection{Quantitative Results}
% Main comparison table + 1--2 figures.
\todo{Results table and figures.}

\subsection{Qualitative Analysis}
% Learned CPTs / states vs. intuition; example inferences; failure cases.
\todo{Interpretation.}

\subsection{Convergence and Scalability}
% EM / Baum--Welch log-likelihood curves; sensitivity to initialization;
% runtime; would more data help?
\todo{Convergence and scalability.}

% Example figure:
% \begin{figure}[t]
%   \centering
%   \includegraphics[width=0.48\linewidth]{em_loglik.png}
%   \caption{EM log-likelihood over iterations for 5 random restarts.}
%   \label{fig:em-convergence}
% \end{figure}
